% Study C (round 5): the compact form. One panel in figure 11's graph grammar: the twelve relabellings of X0 that
% are rotations, as the twelve nodes of the T graph (four g-triangles, solid with heads; the half-turn c = (12)(34)
% as dotted edges), the methane drawn at e and at g (same position, digits moved), the loop as one g-edge; beside
% it the mode discs on the same twelve nodes for one member of each species (A, E cos, T0), and the bundles.
\documentclass[10pt,border=1mm]{standalone}
\usepackage[T1]{fontenc}
\usepackage{figurestyle}
\input{primitives.tex}
\input{molecules.tex}
\input{numbers.tex}
\def\modeA{0/0/1.0000,1/0/1.0000,2/0/1.0000,1/1/1.0000,0/1/1.0000,2/1/1.0000,0/2/1.0000,1/2/1.0000,2/2/1.0000,3/0/1.0000,3/1/1.0000,3/2/1.0000}
\def\modeEcos{0/0/1.0000,1/0/-0.5000,2/0/-0.5000,1/1/1.0000,0/1/-0.5000,2/1/-0.5000,0/2/-0.5000,1/2/-0.5000,2/2/1.0000,3/0/-0.5000,3/1/-0.5000,3/2/1.0000}
\def\modeEsin{0/0/0.0000,1/0/-0.8660,2/0/0.8660,1/1/0.0000,0/1/0.8660,2/1/-0.8660,0/2/-0.8660,1/2/0.8660,2/2/0.0000,3/0/0.8660,3/1/-0.8660,3/2/0.0000}
\def\modeTa{0/0/1.0000,1/0/-0.0000,2/0/0.0000,1/1/-1.0000,0/1/0.0000,2/1/-0.0000,0/2/-0.0000,1/2/0.0000,2/2/-1.0000,3/0/-0.0000,3/1/0.0000,3/2/1.0000}
\def\modeTb{0/0/-0.0000,1/0/0.0000,2/0/-1.0000,1/1/0.0000,0/1/1.0000,2/1/-0.0000,0/2/0.0000,1/2/-1.0000,2/2/-0.0000,3/0/1.0000,3/1/-0.0000,3/2/0.0000}
\def\modeTc{0/0/0.0000,1/0/-1.0000,2/0/-0.0000,1/1/0.0000,0/1/-0.0000,2/1/-1.0000,0/2/1.0000,1/2/0.0000,2/2/-0.0000,3/0/0.0000,3/1/1.0000,3/2/-0.0000}
\def\twelveLabels{0/0/{1234},1/0/{1342},2/0/{1423},1/1/{2143},0/1/{2314},2/1/{2431},0/2/{3124},1/2/{3241},2/2/{3412},3/0/{4132},3/1/{4213},3/2/{4321}}

\newcommand\sitepic[3]{\begin{scope}[shift={(#1,#2)}]
  \foreach \t/\cx/\cy in {0/0/0,1/1.15/0,2/0/-1.05,3/1.15/-1.05} {\draw[fine] ($(\cx,\cy)+(90:.36)$)--($(\cx,\cy)+(210:.36)$)--($(\cx,\cy)+(330:.36)$)--cycle;}
  \foreach \t/\k/\a in #3 {\pgfmathsetmacro{\cx}{ifthenelse(\t==1||\t==3,1.15,0)}\pgfmathsetmacro{\cy}{ifthenelse(\t>=2,-1.05,0)}\pgfmathsetmacro{\ang}{90+120*\k}
    \pgfmathsetmacro{\r}{.42*.36*sqrt(abs(\a))}
    \ifdim\r cm>0.03cm \ifdim\a pt>0pt \fill[PlateInk] ($(\cx,\cy)+(\ang:.36)$) circle[radius=\r]; \else \draw[fine,fill=white] ($(\cx,\cy)+(\ang:.36)$) circle[radius=\r]; \fi
    \else \draw[line width=.5pt,draw=PlateInk] ($(\cx,\cy)+(\ang:.28)$)--($(\cx,\cy)+(\ang:.44)$); \fi}
  \end{scope}}
\begin{document}
\begin{tikzpicture}[plate,jn/.style={line width=.5pt,draw=PlateInk}]
\path[use as bounding box] (0,0) rectangle (15,8.2);
% the graph of T: four g-triangles (nodes = relabellings), the molecule at e and at g
\begin{scope}[shift={(2.8,4.9)}]
  \foreach \t/\cx/\cy in {0/0/0,1/2.3/0,2/0/-2.2,3/2.3/-2.2} {
    \foreach \k/\n in {0/1,1/2,2/0} {\draw[tpath,shorten <=7pt,shorten >=7pt] ($(\cx,\cy)+({90+120*\k}:.8)$) to[bend right=18] ($(\cx,\cy)+({90+120*\n}:.8)$);}}
  \foreach \t/\k/\l in \twelveLabels {\pgfmathsetmacro{\cx}{ifthenelse(\t==1||\t==3,2.3,0)}\pgfmathsetmacro{\cy}{ifthenelse(\t>=2,-2.2,0)}
    \node[circle,draw=PlateInk,fill=white,line width=.5pt,inner sep=0pt,minimum size=5mm,font=\fontsize{6}{7}\selectfont] at ($(\cx,\cy)+({90+120*\k}:.8)$) {\l};}
  \node[sub,anchor=north] at (1.15,-3.3) {the twelve rotations of $T$ as relabellings of $\X_0$; solid edges $g$ (a third-turn), each triangle a $g$-orbit};
  \node[sub,anchor=south] at (1.15,1.05) {the loop $\ell_g$ is one solid edge: from $e=1234$ to $g=2314$};
\end{scope}
\begin{scope}\molsolid\mollabelsinside\pic[scale=.5] at (7.3,6.2) {mol3d-ch4-ref};\end{scope} \node[sub,anchor=north] at (7.3,5.3) {at $e$};
\begin{scope}\molsolid\mollabelsinside\def\molA{2}\def\molB{3}\def\molC{1}\def\molD{4}\pic[scale=.5] at (7.3,3.3) {mol3d-ch4-ref};\end{scope} \node[sub,anchor=north] at (7.3,2.4) {at $g$: same position, digits moved};
% modes on the same twelve nodes
\node[sub,anchor=north] at (11.6,8.1) {modes on the twelve nodes, one member of each species};
\sitepic{9.3}{7.2}{\modeA} \node[sub,anchor=north] at (9.9,5.65) {$\mathrm A$};
\sitepic{11.6}{7.2}{\modeEcos} \node[sub,anchor=north] at (12.2,5.65) {$\mathrm E$};
\sitepic{13.5}{7.2}{\modeTa} \node[sub,anchor=north] at (14.1,5.65) {$\mathrm T$};
\node[sub,anchor=north,align=center] at (11.9,5.2) {along $\ell_g$ the pattern moves one step round each triangle:\\$\mathrm A$ unchanged, $\mathrm E$ by a third ($\omega$), the $\mathrm T$ member into the next};
% bundles, compact
\def\yA{3.4}\def\yE{2.5}\def\yT{1.4}
\draw[heavy] (9.3,\yA)--(10.3,\yA); \node[lab,anchor=base east] at (9.15,{\yA-.1}) {$\mathrm A$}; \node[sub,anchor=base west] at (10.45,{\yA-.1}) {$5$: $(\mathrm A,\mathrm A)$};
\draw[heavy] (9.3,\yE)--(10.3,\yE) (9.3,{\yE+.09})--(10.3,{\yE+.09}); \node[lab,anchor=base east] at (9.15,{\yE-.055}) {$\mathrm E$}; \node[sub,anchor=base west] at (10.45,{\yE-.055}) {$1+1$: $({}^1\mathrm E,{}^2\mathrm E)$, $({}^2\mathrm E,{}^1\mathrm E)$};
\draw[heavy] (9.3,\yT)--(10.3,\yT) (9.3,{\yT+.09})--(10.3,{\yT+.09}) (9.3,{\yT+.18})--(10.3,{\yT+.18}); \node[lab,anchor=base east] at (9.15,{\yT-.01}) {$\mathrm T$}; \node[sub,anchor=base west] at (10.45,{\yT-.01}) {$3$: $(\mathrm T,\mathrm T)$, $\mathfrak d=3$, entangled};
\node[sub,anchor=north] at (11.9,.9) {spin states per level (figure 8's rule) and their isomer names};
\end{tikzpicture}
\end{document}
